\chapter{Randomized Global Min Cut}

\begin{problem}[Global Min Cut]
Given an undirected graph $G = (V,E,c)$, find a cut $(X, V\setminus X)$ that minimizes the total weight of edges crossing the cut, i.e., $\sum_{e \in E(X, V\setminus X)} c(e)$.
\end{problem}

With algorithms of $s$-$t$ min cut as a subroutine, we can solve the global min cut problem by running $s$-$t$ min cut for a fixed $s$ and all $t \in V\setminus \{s\}$, and taking the minimum among these cuts. This approach has a time complexity of $O(n)$ times the time complexity of the $s$-$t$ min cut algorithm, i.e. for general graphs, $n\cdot O(\min\{Cmn, mn^2\})$, and for unit capacity graphs, $n\cdot O(\min\{m^{3/2}, mn^{2/3}\})$. The implementations and proofs are done in Chapter~\ref{chap:maxflow}.

However, we hope to solve global min cut without specifying some $s$ and $t$.

Here we focus on unit capacity graphs.
As shown in homework,
\begin{lemma}
	Let $(X, V\setminus X)$ be one of the global min cuts.
	The number of edges that cross the cut $(X, V\setminus X)$ is at most $2m/n$.
\end{lemma}
\begin{proof}
	\[
		\text{number of edges crossing the cut} \leq \min_{u\in V}\{\deg(u)\}
		\leq \text{average degree} = \frac{2m}{n}.
	\]
\end{proof}

\section{Karger's Random Contraction Algorithm}

Random contraction is a simple randomized algorithm for finding a global min cut. The algorithm works as follows:
\begin{itemize}
	\item pick a random edge
	\item contract it
	\item repeat until only two vertices remain
	\item the remaining two vertices represent a cut, and we return it.
\end{itemize}

We're good as long as we never contract an edge that crosses the global min cut. If we fix the cut $(X, V\setminus X)$, the probability that the first contraction is in the cut is at most $\frac{2m/n}{m} = \frac{2}{n}$. This is independent from the number of edges. As one contraction reduces the number of vertices by 1, we can safely repeat the calculation for all contractions:

\begin{align*}
	\Pr(\text{all contractions are safe}) & \geq \left(1-\frac{2}{n}\right)\left(1-\frac{2}{n-1}\right)\cdots \left(1-\frac{2}{3}\right) \\
	                                      & = \frac{n-2}{n}\cdot \frac{n-3}{n-1}\cdots \frac{1}{3}                                       \\
	                                      & = \frac{2}{n(n-1)}.
\end{align*}

Therefore, the probability that we find the global min cut in $k$ executions is
\[
	1 - (1 - 2/n(n-1))^k \approx 1 - (1-2/n^2)^k \approx 1 - e^{-2k/n^2}.
\]

Set $k = 5n^2\log n$, the probability of success is at least $1 - e^{-10\log n} = 1 - n^{-10}$.

Every contraction can be implemented in $O(n)$ time using Union Find data structure, and there are $n-2$ contractions, so the time complexity of one execution is $O(n^2)$. These together give us the following result:

\begin{theorem}
	With high probability, Karger's random contraction algorithm finds the global min cut in $O(n^4\log n)$ time.
\end{theorem}

\section{Karger's and Stein's Algorithm}

$O(n^4\log n)$ is not good enough. We can do better by modifying Karger's algorithm as follows:

\begin{algorithm}[H]
	\caption{Karger's and Stein's Algorithm}
	\If{$n \leq 10$}{
		\Return Min Cut of $G$ using any algorithm\;
	}
	$G' \gets G$\;
	$G_1 \gets$ random contraction of $G'$ until $|V(G_1)| = \lceil n/\sqrt{2} + 1 \rceil$\;
	$c_1 \gets \text{KargerStein} (G_1)$\;
	$G_2 \gets$ random contraction of $G'$ until $|V(G_2)| = \lceil n/\sqrt{2} + 1 \rceil$\;
	$c_2 \gets \text{KargerStein} (G_2)$\;
	\Return $\min\{c_1, c_2\}$\;
\end{algorithm}

Time complexity is easily to analyze. Let $T(n)$ be the time complexity of Karger's and Stein's algorithm on a graph with $n$ vertices. Then we have

\[ T(n) = 2T(\frac{n}{\sqrt{2}}) + O(n^2) \implies T(n) = O(n^2 \log n). \]

However, the success probability of one trial is harder. To solve this, we need to look at a short diversion.

\section{Galton and Watson's Problem}

Galton and Watson's problem is a classic problem in probability theory that models the extinction of family names. The problem can be stated as follows:

\begin{problem}[Galton and Watson's Problem]
Generation 0 starts with a single individual. Each individual in generation $i$ produces $\zeta\in \{0,1,2,\ldots\}$ offspring in generation $i+1$, where $\zeta$ is a random variable with a given distribution. The process continues indefinitely. What is the probability that the family name eventually becomes extinct?
\end{problem}

Denote $X_i$ as the number of individuals in generation $i$. We state here, without proof, the following result:

\begin{theorem}[Galton and Watson's Theorem] \label{thm:gw}
	Denote $\mu = \E[\zeta]$ as the expected number of offspring per individual. Then eventual extinction ($X_i = 0$ for some $i$) occurs
	\begin{itemize}
		\item with probability 1 if $\mu < 1$;
		\item with probability 1 if $\mu = 1$ unless $\Pr(\zeta = 1) = 1$ (in which case it's 0);
		\item with probability $1-c$ if $\mu > 1$, where $c$ depends on the distribution of $\zeta$.
	\end{itemize}
\end{theorem}

In Karger-Stein Algorithm, we stop at $\lceil n/\sqrt{2} + 1 \rceil$ vertices. The probability that the global min cut is not contracted is at least

\[
	\left(1-\frac{2}{n}\right)\left(1-\frac{2}{n-1}\right)\cdots \left(1-\frac{2}{\lceil n/\sqrt{2} + 1 \rceil + 1}\right) \approx \frac{1}{2}.
\]

We can view the contractions as a branching process, where each branch corresponds to a successful contraction that does not contract the global min cut.
We independently run two branches with probability $\frac{1}{2}$ of success. Therefore, the number of `offspings' follow the distribution of $\Pr(\zeta = 0) = \Pr(\zeta = 2) = 1/4, \Pr(\zeta = 1) = 1/2$. The expected number of offsprings is $\mu = 1$, and $\Pr(\zeta = 1) \neq 1$, so by Theorem~\ref{thm:gw}, the branching will eventually go extinct.

However, we only care about the probability that it goes extinct in the first $2\log n$ generations. Denote $P_i$ as the probability that the branching process still exists after $i$ generations. We have

\begin{align*}
	P_d & = 1-(1-\frac{P_{d-1}}{2})^2     \\
	    & = P_{d-1} - \frac{P_{d-1}^2}{4} \\
\end{align*}

\begin{lemma}
	\[ P_d = \Theta(\frac{1}{d}) .\]
\end{lemma}
\begin{proof}
	We prove by induction, that $P_d \geq \frac{1}{d+1}$.

	$P_0 = 1\geq 1$. Assume $P_{d-1} \geq \frac{1}{d}$.

	As $x - x^2/4$ is increasing for $x\in [0,1]$, we have

	\begin{align*}
		P_d & = P_{d-1} - \frac{P_{d-1}^2}{4}     \\
		    & \geq \frac{1}{d} - \frac{1}{4d^2}   \\
		    & \geq \frac{1}{d} - \frac{1}{d(d+1)} \\
		    & \geq \frac{1}{d+1}.
	\end{align*}
\end{proof}

Therefore, for a $O(n^2\log n)$ time algorithm, the probability of success is at least $P_{2\log n} = \Theta(1/\log n)$. We can repeat the algorithm for $O(\log^2 n)$ times to boost the probability of success to $1 - 1/\text{poly}(n)$.

\begin{theorem}
	With high probability, KargerStein algorithm finds the global min cut in $O(n^2\log^3 n)$ time.
\end{theorem}
